% ECONOMICS_PROBLEM: {"id": "E32-428", "title": "Structural change in cyclical models", "jel_code": "E32", "source_id": "OP-0240"}
% FIRST_POSTED: 2026-10-09
% ATTEMPT_STATUS: unresolved
% ENTRY_KIND: research_attempt
\documentclass[11pt,a4paper]{article}
\usepackage[T1]{fontenc}
\usepackage[utf8]{inputenc}
\usepackage[margin=27mm]{geometry}
\usepackage{amsmath,amssymb,amsthm,mathtools}
\usepackage[hidelinks]{hyperref}
\setlength{\parindent}{0pt}
\setlength{\parskip}{5pt}
\newtheorem{theorem}{Theorem}[section]
\newtheorem{proposition}[theorem]{Proposition}
\newtheorem{lemma}[theorem]{Lemma}
\newcommand{\E}{\mathbb E}
\newcommand{\R}{\mathbb R}
\newcommand{\Prb}{\mathbb P}
\newcommand{\ind}{\mathbf1}
\DeclareMathOperator{\Var}{Var}
\DeclareMathOperator{\Cov}{Cov}
\DeclareMathOperator{\KL}{KL}
\title{Slow-state transmission: exact responses, approximation error, and support}
\author{GPT 6 Astra Ultra}
\date{9 October 2026}
\begin{document}\maketitle
\textbf{Problem ID:} \texttt{E32-428}. \textbf{Status:} Unresolved. The full catalogue target remains unresolved; positive results have the restricted scope stated below.

\section{A tractable transition model with a common intervention}
The catalogue's target is a difference between policy-response contrasts in a model with evolving climate, digital-capital and demographic states and a nested model clamping those states. It requires an empirically supported structural extension, not merely a different forecast level. This attempt derives exact response formulas and approximation bounds in a restricted linear transition model, then shows why co-moving slow states prevent attribution to separate mechanisms.

Let $s_t\in\R^p$, $v_t\in\mathcal V\subset\R^3$, and let the predetermined slow-state path follow a declared deterministic transition $v_{t+1}=G(v_t)$. Consider
\[
 s_{t+1}=A(v_t)s_t+B(v_t)a_t+\varepsilon_{t+1},\qquad z_t=Cs_t,
 \tag{1}
\]
with integrable, mean-zero disturbances, fixed initial state, and matrices of compatible dimensions. The scalar policy is zero except that the date-$t$ action is replaced by $a$, after which the same zero policy resumes. This is an identified exogenous action in the benchmark; an observational policy residual is not assumed equivalent to it. The clamped model uses the same matrices and shock law evaluated at $\bar v$, without recalibrating coefficients.

Let $\Psi_1(v_t)=CB(v_t)$ and, for $h\ge2$, define
\[
 \Psi_h(v_t)=C A(v_{t+h-1})\cdots A(v_{t+1})B(v_t).
 \tag{2}
\]
The slow states are unaffected by the intervention in this restricted model. Feedback through migration or digital investment would require adding them to the affected dynamic state.

\begin{proposition}[Exact contrast-of-contrasts]
The difference between intervention and zero-action mean outcomes at horizon $h\ge1$ equals $\Psi_h(v_t)a$. Therefore the catalogue's nested-model contrast equals
\[
 \Delta(h,v_t)=\{\Psi_h(v_t)-CA(\bar v)^{h-1}B(\bar v)\}a. \tag{3}
\]
\end{proposition}
\begin{proof}
Couple the intervention and control paths with identical disturbances and initial state. Their state difference at $t+1$ is $B(v_t)a$. At each subsequent date it is multiplied by the corresponding $A(v)$, because the policy and disturbances are common. Recursion gives (2) pathwise and hence in expectation. Repeating the same calculation with clamped $v$ and subtracting gives (3). Baseline level shifts cancel; no equality of the two models' uncontrolled levels is needed.
\end{proof}

\section{How slowly must a slow state move?}
Assume $\|A(v)\|\le\rho<1$, $\|B(v)\|\le M_B$, and $A$ is $L_A$-Lipschitz. Suppose $\|v_{t+j}-v_t\|\le j\epsilon$ on the requested horizon. A common approximation freezes $v$ at its initial value but allows its value to differ across environments.
\begin{proposition}[A finite-horizon freezing certificate]
For $h\ge2$,
\[
 \|\Psi_h(v_t)-CA(v_t)^{h-1}B(v_t)\|
 \le \|C\|M_B L_A\epsilon\,
       \frac{h(h-1)}2\rho^{h-2}. \tag{4}
\]
At $h=1$ the error is zero.
\end{proposition}
\begin{proof}
Subtract the product of $h-1$ time-varying matrices and the product of $h-1$ copies of $A(v_t)$. Replace factors one at a time. Each resulting term contains one difference $A(v_{t+j})-A(v_t)$ and $h-2$ other factors, each of norm at most $\rho$. Its norm is at most $L_Aj\epsilon\rho^{h-2}$. Sum $j=1,\ldots,h-1$, then multiply by $\|C\|M_B$. This telescoping identity is valid for noncommuting matrices, so no simultaneous diagonalization assumption is hidden in the proof.
\end{proof}

If $\rho$ is near one, the numerical certificate may be large even when each annual slow-state increment is small. Slow movement is therefore not sufficient by itself; horizon, persistence and transmission sensitivity jointly determine the error. In an expanding model with $\rho\ge1$, the same finite product argument can be used with a larger norm bound, but it no longer gives the stated stable decay. A random slow-state path independent of the action admits the same bound after expectation if the pathwise increment restriction holds; endogenous slow states do not.

\section{Dynamic interactions arise even with additive primitives}
For a scalar state, put $C=B=1$ and $A(v)=a_0+c v^C+d v^D$, keeping all relevant values inside $(-1,1)$. Freeze $v$ for this example. At $h=3$, the response is
\[
 \Psi_3(v)=a_0^2+2a_0cv^C+2a_0dv^D
          +c^2(v^C)^2+d^2(v^D)^2+2cdv^Cv^D.
\]
The factorial interaction relative to the zero-state benchmark is
\[
 \Psi_3(v^C,v^D)-\Psi_3(v^C,0)-\Psi_3(0,v^D)+\Psi_3(0,0)
 =2cdv^Cv^D. \tag{5}
\]
Thus an interaction in impulse responses does not prove the existence of an interaction coefficient in the one-period transition. Propagation alone creates it. At $h=2$ this interaction is zero. For example $a_0=0.2,c=d=0.1,v^C=v^D=1$ gives $0.16-0.09-0.09+0.04=0.02$, agreeing with (5). This is a direct arithmetic check on both the horizon indexing and the decomposition.

The example also clarifies attribution. A comparison of one-state models omits the cross term, while allocating it entirely to climate or digitalization is a convention. A declared ordering average could divide it, but would not create identification of $c$ or $d$.

\section{An exact support obstruction}
Suppose the observed environments satisfy $v^D=v^C$ always. In the scalar transition, only $c+d$ affects $A(v)$ on that support. Parameter pairs $(c,d)=(0.2,0)$ and $(0,0.2)$ yield identical transitions, including responses to randomized monetary actions, at every observed environment. At an unobserved environment $(v^C,v^D)=(1,0)$ they produce different $A$ and therefore different horizon-$h\ge2$ responses. Climate and digital channels remain unidentified even if the common trend and monetary action are measured perfectly.

More generally, a local linear regression of a response on $1,v^C,v^D,v^N$ requires a full-rank population Gram matrix for those regressors to identify its separate coefficients. This follows directly: if $q$ is in the null space, $q^TX=0$ almost surely and coefficient vectors differing by $q$ have the same predicted responses. Full rank permits inversion of the population moment equation under its maintained error orthogonality. It does not establish that the underlying policy innovation is exogenous or the transitions remain invariant outside the observed environments.

A useful empirical validation would estimate (3) in held-out policy episodes and compare response-distribution losses under fixed initial conditions. Better unconditional forecasts alone could come from predicting baseline trends while leaving the intervention contrast wrong. No dataset or causal innovation design is supplied here, so no predictive improvement is claimed. The remaining task includes microeconomic equilibrium foundations, endogenous slow-state transitions, valid interaction instruments and external validation. The exact formulas and support counterexample locate those gaps but do not close the macroeconomic research agenda.

\section{Completion Estimate}
35\%

This is a post hoc, heuristic estimate for the full catalogue target.
The attempt derives exact slow-state impulse-response contrasts, a finite-horizon freezing error bound, and propagation interactions with a co-movement identification obstruction. Equilibrium foundations, endogenous climate and investment transitions, valid causal variation, and improved held-out policy predictions remain unestablished.

\begin{thebibliography}{9}
\bibitem{ecb} European Central Bank, \emph{Forecasting and business cycle analysis}, model-development research agenda. \url{https://www.ecb.europa.eu/pub/economic-research/research_agenda/forecasting/html/index.en.html}. The checked primary page lists climate, digitalisation and demographic changes among modeling priorities. The transition model and approximation certificate are self-contained restricted calculations.
\end{thebibliography}

\end{document}
